\chap{物理约束温度场世界模型}

\sect{状态转移表示}

离散温度场定义为
\begin{equation}
  \bm{s}_n=[T_1^n,T_2^n,\ldots,T_N^n]^{\mathsf T},\qquad N=41.
\end{equation}
动作 $a_n$ 为下一时间步的环境温度。参数向量包含边界、物性尺度、厚度和时间步长。网络预测温度增量
\begin{equation}
  \Delta\widehat{\bm{T}}_n=f_{\theta}(\bm{s}_n,a_n,\bm{p}),\qquad
  \widehat{\bm{s}}_{n+1}=\bm{s}_n+\Delta\widehat{\bm{T}}_n.
  \label{eq:residual_wm}
\end{equation}
增量形式保留了相邻状态的连续性，也使网络输出尺度与单步热响应一致。

主模型采用三隐藏层多层感知机，每层 128 个单元，激活函数为 SiLU。这里的“残差”指网络预测状态增量，网络内部没有残差块。$h,\varepsilon$ 参数化的输入维数为 $41+1+7=49$，$\Heff$ 参数化为 $41+1+6=48$。输入状态、动作和参数依据训练集统计量标准化，输出温度增量按节点标准化。当前固定网格只有 41 个节点，MLP 可以直接处理完整状态，并适合在控制器内批量评价候选动作。

\begin{figure}[H]
  \centering
  \begin{tikzpicture}[
    node distance=9mm and 12mm,
    box/.style={draw=accent,thick,rounded corners=2pt,minimum height=10mm,minimum width=29mm,align=center,fill=accent!5},
    arrow/.style={-{Latex[length=2.2mm]},thick,draw=accent}
  ]
    \node[box] (state) {温度场 $\bm{s}_n$};
    \node[box,below=of state] (action) {炉温动作 $a_n$};
    \node[box,below=of action] (param) {物性与 $\Heff$};
    \node[box,right=18mm of action,minimum width=39mm] (model) {残差世界模型\\$f_\theta$};
    \node[box,right=18mm of model] (next) {预测场 $\widehat{\bm{s}}_{n+1}$};
    \node[box,above=of next] (loss) {数据与物理损失};
    \draw[arrow] (state.east) -- (model.west);
    \draw[arrow] (action.east) -- (model.west);
    \draw[arrow] (param.east) -- (model.west);
    \draw[arrow] (model) -- (next);
    \draw[arrow] (next) -- (loss);
    \draw[arrow] (next.south) |- ++(0,-8mm) -| (state.south);
  \end{tikzpicture}
  \caption{受控递归温度场世界模型}
  \label{fig:architecture}
\end{figure}

\sect{闭环多步训练}

训练样本以长度为 $K$ 的连续片段组织。第一步输入真实状态，后续输入采用网络上一时刻输出：
\begin{equation}
  \widehat{\bm{s}}_{n+j+1}=F_{\theta}(\widehat{\bm{s}}_{n+j},a_{n+j},\bm{p}_{n+j}),
  \quad j=0,\ldots,K-1,
\end{equation}
其中 $\widehat{\bm{s}}_n=\bm{s}_n$，本文取 $K=5$。令 $\sigma_{\Delta T,i}$ 为训练集中节点 $i$ 的单步温度增量标准差。多步数据损失为
\begin{equation}
  \mathcal{L}_{\mathrm{data}}
  =\frac{1}{KN}\sum_{j=1}^{K}\sum_{i=1}^{N}
  \left(
    \frac{\widehat{T}_{i}^{n+j}-T_i^{n+j}}{\sigma_{\Delta T,i}}
  \right)^2.
  \label{eq:data_loss}
\end{equation}
这项损失直接约束部署时会出现的模型生成状态。检查点依据验证集 300 步完整轨迹 RMSE 选择。该选择标准与工艺推演目标保持一致。

\begin{table}[H]
  \centering
  \caption{世界模型的正式训练设置}
  \label{tab:training_config}
  \begin{tabular}{ll@{\hspace{1.8em}}ll}
    \toprule
    项目 & 设置 & 项目 & 设置 \\
    \midrule
    隐藏层 & $3\times128$ & 激活函数 & SiLU \\
    优化器 & AdamW & 初始学习率 & $10^{-3}$ \\
    权重衰减 & $10^{-6}$ & 批量大小 & 512 \\
    训练轮数 & 120 & 展开长度 & $K=5$ \\
    梯度范数上限 & 5 & 随机种子 & 42、7、123 \\
    学习率调度 & 平台期减半 & 最低学习率 & $10^{-5}$ \\
    \bottomrule
  \end{tabular}
\end{table}

\sect{离散物理约束}

内部节点物理残差采用与数值生成器一致的温度增量形式。对预测状态 $\widehat{\bm{s}}_{n+1}$，定义
\begin{equation}
  R_i=\widehat{T}_i^{n+1}-T_i^n
  -\frac{\Delta t}{\rho c_{p,i}^{n}\Delta x^2}\left[
  k_{i+1/2}^{n}(\widehat{T}_{i+1}^{n+1}-\widehat{T}_{i}^{n+1})
  +k_{i-1/2}^{n}(\widehat{T}_{i-1}^{n+1}-\widehat{T}_{i}^{n+1})
  \right].
  \label{eq:temperature_residual}
\end{equation}
$R_i$ 的量纲为温度。表面节点残差由式\eqref{eq:surface_discrete_balance}左、右两端之差构成，并使用对应的对流辐射热流。残差按 $\sigma_{\Delta T,i}$ 归一化，记为 $\widetilde{R}_i=R_i/\sigma_{\Delta T,i}$。多步物理损失写为
\begin{equation}
  \mathcal{L}_{\mathrm{phy}}
  =\frac{1}{K}\sum_{j=1}^{K}
  \frac{1}{N}\sum_{i=1}^{N}\widetilde{R}_{i,n+j}^{2}.
\end{equation}
总损失为
\begin{equation}
  \mathcal{L}=\mathcal{L}_{\mathrm{data}}+\lambda_{\mathrm{phy}}\mathcal{L}_{\mathrm{phy}}.
  \label{eq:total_loss}
\end{equation}
单种子预扫描包含 $0$、$10^{-3}$、$10^{-2}$ 和 $10^{-1}$。$10^{-1}$ 的场误差已明显增大，所以三随机种子重复实验比较 $0$、$10^{-3}$ 和 $10^{-2}$。验证轨迹 RMSE 的均值选择 $\lambda_{\mathrm{phy}}=10^{-3}$ 作为正式设置。

\sect{动态边界的可辨识参数化}

\noindent\textbf{命题 1（单步边界转移的局部不可辨识性）}
在固定 $T_s$、$\Tinfty$ 且 $h,\varepsilon$ 均未知时，边界热流只依赖标量组合 $\Heff$，映射 $(h,\varepsilon)\mapsto q_s$ 非单射。

由式\eqref{eq:heff_flux}定义
\begin{equation}
  \Heff(t)=h(t)+\varepsilon(t)h_{\mathrm{rad}}(T_s(t),\Tinfty(t)).
  \label{eq:heff}
\end{equation}
在给定 $T_s$ 和 $\Tinfty$ 的瞬时转移中，边界热流对两个参数的灵敏度为
\begin{equation}
  \frac{\partial q_s}{\partial(h,\varepsilon)}
  =(\Tinfty-T_s)\begin{bmatrix}1 & h_{\mathrm{rad}}\end{bmatrix}.
  \label{eq:sensitivity}
\end{equation}
当 $\Tinfty\ne T_s$ 时，该行向量的秩为 1。温度响应只能确定一个标量组合。任取扰动 $\delta\varepsilon$，只要
\begin{equation}
  \delta h=-h_{\mathrm{rad}}\delta\varepsilon,
\end{equation}
就能在当前状态保持相同热流，命题得证。逐时自由变化的 $h(t)$ 和 $\varepsilon(t)$ 因而没有唯一分解。该结论限定于本文的均匀灰体边界、已知表面温度和已知炉温。若参数在一段时间内保持常数，变化的 $h_{\mathrm{rad}}$ 可提供额外历史信息。

这里需要区分两种充分性。对已执行动作下的前向状态转移，$\Heff$ 是当前边界热流的充分输入。对候选炉温的反事实规划，$h_{\mathrm{rad}}$ 随炉温改变，当前 $\Heff$ 不能单独确定新动作下的边界热流。因此，$\Heff$ 不构成整个控制问题的充分统计量。

基于该结论，本文构建第二类世界模型
\begin{equation}
  \widehat{\bm{s}}_{n+1}
  =F_{\theta}^{H}(\bm{s}_n,a_n,\Heff_n,s_k,s_c,L,\Delta t).
  \label{eq:heff_model}
\end{equation}
该表示消除了训练输入中的等价方向。网络无需学习同一边界热流对应的多组参数组合。闭环规划阶段使用增广状态滤波器保持 $h$ 与 $\varepsilon$ 的后验集合，并由每个集合成员计算候选动作下的 $\Heff(a')$。

\sect{评价指标}

完整轨迹滚动 RMSE 定义为
\begin{equation}
  \mathrm{RMSE}_{\mathrm{roll}}
  =\sqrt{\frac{1}{MN\tau}\sum_{m=1}^{M}\sum_{n=1}^{\tau}
  \left\|\widehat{\bm{s}}_{m,n}-\bm{s}_{m,n}\right\|_2^2}.
\end{equation}
同时报告逐轨迹 RMSE 的中位数和 95\% 分位数、最大绝对误差、最终状态 RMSE 与离散物理残差。最大值原理指标采用全局包络。下一状态若超出当前全场温度与环境温度共同确定的上下界，则记为一次违反。该定义允许节点因邻域导热合法升温或降温。

正式模型结果使用 42、7 和 123 三个训练种子。参数比较在相同测试轨迹上进行，表中的“均值$\pm$标准差”默认表示三个训练种子的均值与样本标准差。模型选择只访问验证集。BDF 参考集和控制环境不参与训练。闭环实验重复使用相同的 BDF 场景比较不同控制器，所以控制差异采用场景配对统计。
